\documentclass[a4paper,12 pt]{report}
\usepackage[T1]{fontenc}
\usepackage[utf8]{inputenc}
\usepackage{lmodern}
\usepackage{listings}
\usepackage{graphicx}
\usepackage{float}
\usepackage{subcaption}
\usepackage{hyperref}
\usepackage{wrapfig}
\usepackage{fancyhdr}

% forza le footnote a stare il più in basso possibile
\usepackage[bottom]{footmisc}


%% STILE LISTINGS

\usepackage{xcolor}

\definecolor{codegreen}{rgb}{0,0.6,0}
\definecolor{codegray}{rgb}{0.5,0.5,0.5}
\definecolor{codepurple}{rgb}{0.58,0,0.82}
\definecolor{backcolour}{rgb}{0.95,0.95,0.92}

\lstdefinestyle{mystyle}{
    backgroundcolor=\color{backcolour},   
    commentstyle=\color{codegreen},
    keywordstyle=\color{magenta},
    numberstyle=\tiny\color{codegray},
    stringstyle=\color{codepurple},
    basicstyle=\ttfamily\footnotesize,
    breakatwhitespace=false,         
    breaklines=true,                 
    captionpos=b,                    
    keepspaces=true,                 
    numbers=left,                    
    numbersep=5pt,                  
    showspaces=false,                
    showstringspaces=false,
    showtabs=false,                  
    tabsize=2
}

\lstset{style=mystyle}

%% -----

% mostra le subsubsection nell'indice
\setcounter{tocdepth}{3}
\setcounter{secnumdepth}{3}

% Resetta la numerazione dei chapter quando
% una nuova part viene creata
\makeatletter
\@addtoreset{chapter}{part}
\makeatother

% Rimuove l'indentazione quando si crea un nuovo paragrafo
\setlength{\parindent}{0pt}

% footer
\pagestyle{fancyplain}
% rimuove la riga nell'header
\fancyhf{} % sets both header and footer to nothing
\renewcommand{\headrulewidth}{0pt}
\fancyfoot[L]{\href{https://github.com/Typing-Monkeys/AppuntiUniversita}{Typing Monkeys}}
\fancyfoot[C]{\emoji{gorilla}}
\fancyfoot[R]{\thepage}

% configurazione emoji
\usepackage{fontspec}
\usepackage{emoji}
\setemojifont{NotoColorEmoji.ttf}[Path=/usr/share/fonts/truetype/noto/]

\begin{document}
% Thesis frontmatter --------------------------------------------

\thispagestyle{empty} %suppress page number

\noindent % just to prevent indentation narrowing the line width for this line
\includegraphics[width=0.15\textwidth]{frontmatter/imgs/logoUniPg.jpg}
\begin{minipage}[b]{0.7\textwidth}
    \centering
    {\Large \textsc{Universit{\`a} di Perugia}}\\
    \vspace{0.4 em}
    {\large Dipartimento di Matematica e Informatica}
    \vspace{0.6 em}
\end{minipage}%
\includegraphics[width=0.15\textwidth]{frontmatter/imgs/logoDMI.jpg}

\vspace{8 em}

\begin{center}
    {\Huge Appunti \textit{Inserisci Nome}}\\
    \vspace{5 em}
    {\Huge \textbf{Testo}}\\

    \vfill

    \rule{380pt}{.4pt}\\
    \vspace{1.2 em}
    \large{Anno Accademico 2021-2022}\\ %% MODIFICARE
    \vspace{.9 em}
    \small{\textit{Last Update: \today{}}}
\end{center}

% ------------------------------------------------------------------

\tableofcontents

\newlength\longest

\clearpage

\thispagestyle{empty}
\null\vfill

\settowidth\longest{\huge\itshape just as his inclination leads him;}
\begin{center}
    \parbox{\longest}{%
        \raggedright{\LARGE\itshape%
            "Oi, con quanto sentimento\\
            defeco sul tuo naso,\\
            così che ti coli sul mento."\par\bigskip
        }
        \raggedleft\Large{Wolfgang Amadeus Mozart}\par%
    }
\end{center}


\vfill\vfill

\clearpage

%% Aggiungere i capitoli qui sotto
\chapter{Esempio 1}

\section{Introduzione}

Lorem ipsum dolor sit amet. A excepturi mollitia eum iste aperiam qui
perspiciatis debitis qui quia voluptas et beatae voluptatem qui animi molestiae
in explicabo placeat. Ut blanditiis voluptas et perferendis labore id voluptatem
animi. Ad suscipit exercitationem ut nesciunt laudantium ut ullam aspernatur in
veniam enim ut galisum obcaecati hic tempora repudiandae sit quidem ullam. 33
corrupti fugit sed facere dolor rem labore molestiae ab veniam nostrum aut
internos omnis.
\section{Arg1}

Argomento 1 con una sotto sezione

\subsection{Sotto Sezione 1}

Io sono una sottosezione !! \emoji{gorilla}



\end{document}